\clearpage
\SourceLeaf{161}\pdfbookmark[0]{Book Two.}{liber-secundus}
\begin{center}
{\Large BOOK TWO.}\par\bigskip
{\large INVESTIGATION OF THE ORBITS OF CELESTIAL BODIES FROM GEOCENTRIC OBSERVATIONS.}\par\bigskip
\rule{3em}{.3pt}\par\bigskip
\pdfbookmark[1]{Section One.}{liber-secundus-sectio-prima}
{\large SECTION ONE.}\par\medskip
\emph{Determination of the orbit from three complete observations.}\par\bigskip
\hypertarget{native-161-article115}{}\pdfbookmark[2]{115.}{article-115}115.
\end{center}
A complete determination of the motion of a celestial body in its orbit requires \emph{seven} elements. This number is reduced by one, however, if the body's mass is either known or neglected. It will scarcely be possible to avoid neglecting it when determining a hitherto entirely unknown orbit, since all quantities of the order of the perturbations must be set aside until the masses on which they depend are known from other sources. In the present investigation, therefore, we neglect the body's mass and reduce the number of elements to six. It is then clear that to determine an unknown orbit we require as many quantities that depend on the elements and are independent of one another. These quantities must come from positions of the celestial body observed from Earth. Since each such position supplies two data, namely longitude and latitude or right ascension and declination, the simplest course is to use \emph{three geocentric positions}, which in general will suffice to determine the six unknown elements. This problem is to be regarded as the \SourceLeaf{162}most important in this work and will therefore be treated with the utmost care in this section.

In the special case, however, in which the orbital plane coincides with the ecliptic and hence all latitudes, both heliocentric and geocentric, vanish by their nature, the three vanishing geocentric latitudes can no longer be regarded as three independent data. The problem would then remain indeterminate, and infinitely many orbits could satisfy the three geocentric positions. In such a case, therefore, four geocentric longitudes must be given to determine the four remaining unknown elements (the inclination of the orbit and the longitude of the node having dropped out). Although on the principle of indiscernibles there is no reason to expect such a case ever to occur in nature, it is easy to suppose that the problem, which becomes absolutely indeterminate for an orbit that coincides exactly with the plane of the ecliptic, must remain virtually indeterminate \emph{for orbits very slightly inclined to the ecliptic}, owing to the limited precision of the observations. In this situation even the slightest errors in the observations can completely upset the determination of the unknowns. To provide for this case as well, we must choose six different data. To this end we shall show in the second section how to determine the unknown orbit from four observations, two of which are complete while the other two lack latitudes or declinations.

Finally, since all our observations are only approximations to the truth owing to the imperfection of our instruments and senses, an orbit based on just the six strictly necessary data may still be liable to considerable errors. To reduce these errors as far as possible and thus attain the greatest possible precision, our only course is to collect as many of the most precise observations as possible and refine the elements so that they agree as closely as possible with all of them, without requiring exact agreement with any particular observations. In the third section we shall show how to obtain such agreement according to the principles of the calculus of probability, as close as possible everywhere even if nowhere exact.

In this way the determination of orbits, insofar as celestial bodies move in them according to \textsc{Kepler}'s laws, will be brought to every degree of perfection one could desire. The final refinement can only be undertaken \SourceLeaf{163}once the perturbations that the other planets introduce into the motion have also been taken into account in the calculation. In the fourth section we shall briefly indicate how these perturbations should be taken into account insofar as it seems relevant to our purpose.
